\documentclass[runningheads]{llncs}

% ---------------------------------------------------------------
% Include basic ECCV package

% TODO REVIEW: Insert your submission number below by replacing '*****'
% TODO FINAL: Comment out the following line for the camera-ready version
\usepackage[review,year=2026,ID=*****]{eccv}
% TODO FINAL: Un-comment the following line for the camera-ready version
%\usepackage{eccv}

% OPTIONAL: Un-comment the following line for a version which is easier to read
% on small portrait-orientation screens (e.g., mobile phones, or beside other windows)
%\usepackage[mobile]{eccv}


% ---------------------------------------------------------------
% Other packages

% Commonly used abbreviations (\eg, \ie, \etc, \cf, \etal, etc.)
\usepackage{eccvabbrv}

% Include other packages here, before hyperref.
\usepackage{graphicx}
\usepackage{booktabs}
\usepackage{amsmath, amssymb}
\usepackage{xcolor}
\usepackage{subcaption}
\usepackage{algorithm}
\usepackage{algpseudocode}
\usepackage{multirow}
\usepackage{array}
\usepackage{makecell}
\usepackage{enumitem}
\usepackage{colortbl}
\usepackage{pifont}
\usepackage{marvosym}
\usepackage{fontawesome5}
\usepackage{float}
\usepackage{microtype}
\usepackage{tikz}
\usetikzlibrary{shapes.geometric, arrows.meta, positioning, calc, backgrounds, fit}

% The "axessibility" package can be found at: https://ctan.org/pkg/axessibility?lang=en
\usepackage[accsupp]{axessibility}  % Improves PDF readability for those with disabilities.


% ---------------------------------------------------------------
% Hyperref package

% It is strongly recommended to use hyperref, especially for the review version.
% Please disable hyperref *only* if you encounter grave issues.
% hyperref with option pagebackref eases the reviewers' job, but should be disabled for the final version.
%
% If you comment hyperref and then uncomment it, you should delete
% main.aux before re-running LaTeX.
% (Or just hit 'q' on the first LaTeX run, let it finish, and you
%  should be clear).

% TODO FINAL: Comment out the following line for the camera-ready version
\usepackage[pagebackref,breaklinks,colorlinks,citecolor=eccvblue]{hyperref}
% TODO FINAL: Un-comment the following line for the camera-ready version
%\usepackage{hyperref}

% Support for ORCID icon
\usepackage{orcidlink}

% Cross-references (must be loaded after hyperref)
\usepackage[nameinlink,capitalize,noabbrev]{cleveref}


% ---------------------------------------------------------------
% Paper-specific macros
% (If these are already defined in your Overleaf project, remove this block.)

% Method name macros
\newcommand{\method}{\textsc{Kodiak}}
\newcommand{\ours}{\textbf{\textsc{Kodiak} (Ours)}}
\newcommand{\oursshort}{\textsc{Kodiak}}

% Convenience
\newcommand{\xmark}{\textcolor{gray}{\ding{55}}}
\newcommand{\cmark}{\checkmark}
\newcommand{\NA}{\textcolor{gray}{--}}
\newcommand{\result}{\textcolor{red}{\textbf{--}}}
\newcommand{\best}[1]{\textbf{#1}}
\newcommand{\second}[1]{\underline{#1}}

% Delta annotations
\definecolor{kodiakgreen}{RGB}{44, 160, 44}
\definecolor{kodiakred}{RGB}{214, 39, 40}
\newcommand{\gain}[1]{\textcolor{kodiakgreen}{\scriptsize(+#1)}}
\newcommand{\drop}[1]{\textcolor{kodiakred}{\scriptsize(-#1)}}

% Regime symbols
\newcommand{\globe}{\footnotesize{\faGlobe}\normalsize}
\newcommand{\basesym}{\globe}
\newcommand{\fsym}{\textsuperscript{\textdagger}}
\newcommand{\csym}{\textsuperscript{\textdaggerdbl}}

% Table helpers (MOCA-style)
\newcommand{\tablestyle}[2]{\setlength{\tabcolsep}{#1}\renewcommand{\arraystretch}{#2}\centering\scriptsize}
\definecolor{baselinecolor}{gray}{.9}
\newcommand{\baseline}[1]{\cellcolor{baselinecolor}{#1}}
\definecolor{oursrow}{RGB}{218, 238, 255}
\newcommand{\rowours}{\rowcolor{oursrow}}

% Colors for algorithm
\definecolor{kodiakblue}{RGB}{31, 119, 180}

% Math operators
\DeclareMathOperator{\softmax}{softmax}
\DeclareMathOperator{\logsoftmax}{log\text{-}softmax}
\DeclareMathOperator{\sinkhorn}{Sinkhorn}
\newcommand{\norm}[1]{\left\lVert #1 \right\rVert}

% Make \paragraph headings bold upright (LLNCS default is bold italic)
\makeatletter
\renewcommand\paragraph{\@startsection{paragraph}{4}{\z@}%
  {-2.5ex\@plus -1ex \@minus -.2ex}%
  {-1em}%
  {\normalfont\normalsize\bfseries}}
\makeatother

% Reduce overfull hboxes
\emergencystretch=1em


\begin{document}

% ---------------------------------------------------------------
\title{Kodiak: Discrete Codebook Distillation for Low-Data Self-Supervised Adaptation}

\titlerunning{Kodiak: Discrete Codebook Distillation}

% TODO FINAL: Replace with your author list.
\author{Anonymous ECCV 2026 Submission}

\authorrunning{Anonymous}

% TODO FINAL: Replace with your institution list.
\institute{Paper ID *****}

\maketitle


% ==========================================================
% ABSTRACT
% ==========================================================
\begin{abstract}
Self-supervised adaptation of foundation models to small specialized domains remains fragile: continuing standard objectives such as DINO-style distillation on limited domain data often \emph{degrades} transfer quality rather than improving it.
We provide evidence that continuous feature targets contribute to this failure---they are high-dimensional, prone to representational drift, and lack categorical structure to anchor learning when data diversity is limited.
To address this, we propose \textbf{\method{}}, a single-stage framework that replaces continuous targets with \emph{discrete codebook assignments}, stabilized by Sinkhorn-Knopp balanced teacher targets.
A shared visual codebook serves as the target space for both masked patch prediction (dense) and cross-view CLS alignment (global), requiring no external tokenizers or multi-stage pipelines while maintaining full compatibility with standard vision transformer architectures.
We study \method{} in the low-data regime across four diverse visual domains, under both from-scratch and foundation-model adaptation settings.
As demonstrated across multiple evaluation protocols, \method{} consistently outperforms continuous baselines: it improves the foundation model on domains where the standard DINO objective degrades it, and substantially outperforms all baselines when training from scratch.
Our analysis further reveals that component importance is regime-dependent: auxiliary losses are critical for from-scratch training, while the discrete codebook alone drives gains when adapting a strong backbone.
\keywords{Self-supervised learning \and Low-data adaptation \and Discrete codebook \and Vision transformers}
\end{abstract}


% ==========================================================
% 1. INTRODUCTION
% ==========================================================
\section{Introduction}
\label{sec:intro}

Foundation models trained with self-supervised learning (SSL) have become the default initialization for visual recognition.
Yet adapting these models to specialized domains with limited data remains an open problem.
When a practitioner has 5K--50K unlabeled domain images (the reality in remote sensing, texture analysis, and fine-grained recognition), the natural strategy is to continue pretraining with the same SSL objective.
We show that this strategy can \emph{backfire}: continuing DINOv3-style self-distillation on DTD (5.6K images) drops fine-tuning accuracy from $71.12\%$ to $62.44\%$, and on Oxford-Pets (7.4K images) from $86.22\%$ to $80.26\%$.
The foundation model \emph{degrades} when adapted with its own training objective (\cref{fig:adaptation_delta}).

\begin{figure}[t]
\centering
\includegraphics[width=\linewidth]{figures/adaptation_delta.pdf}
\caption{\textbf{When continued pretraining helps vs.\ hurts.}
Fine-tuning accuracy change (pp) relative to the DINOv3 foundation model after continued pretraining on each target domain.
Error bars: std over 3 seeds.
DINOv3 \emph{degrades} DTD ($-$8.7) and Oxford-Pets ($-$6.0); \method{} consistently \emph{improves} the foundation.}
\label{fig:adaptation_delta}
\end{figure}

Why does continuation hurt?
We hypothesize that continuous feature targets are a poor match for the low-data regime.
The student must regress high-dimensional feature vectors from a slowly drifting teacher, with insufficient data diversity to constrain the optimization.
This leads to representational drift and overfitting to surface statistics of the small dataset.
Crucially, this degradation is \emph{not} an artifact of poor hyperparameter tuning: all methods share the same backbone, augmentation pipeline, optimizer, and training budget, and DINOv3 hyperparameters were grid-searched on a per-dataset basis (\cref{sec:setup}).
While other factors contribute (\eg masking strategy, regularization), our controlled study suggests the \emph{target space} is a primary driver of adaptation success in this regime.

This observation motivates a simple design principle: \textbf{use discrete, balanced codebook assignments as the target space for SSL}.
Discrete targets provide three structural advantages over continuous regression: (i)~an information bottleneck that forces categorical decisions rather than high-dimensional matching; (ii)~Sinkhorn balancing that prevents codebook collapse and ensures diverse representation usage; and (iii)~a shared discrete space that can supervise both dense (patch-level) and global (image-level) learning signals.

We instantiate this principle in \textbf{\method{}} (\textbf{KO}debook \textbf{DI}stillation for \textbf{A}daptation; \textit{Kodiak} for short), a single-stage framework built on the standard DINO teacher-student protocol.
The student predicts Sinkhorn-balanced codebook distributions at two levels: (i)~\emph{patch-level}, where masked patches are reconstructed as codebook assignments from teacher targets, and (ii)~\emph{image-level}, where CLS tokens from global and local crops are aligned in the same codebook space.
We evaluate on four domains (5K--60K images) under \emph{continued pretraining} from DINOv3 (the practical adaptation scenario) and \emph{from-scratch training} (a stress test isolating the objective from initialization).

Our contributions are threefold.
\textbf{(1)}~We identify the target space as a primary determinant of adaptation success under limited data, providing controlled evidence that discrete codebook assignments outperform continuous features across four domains, two training regimes, and four evaluation protocols.
\textbf{(2)}~We introduce \method{}, a single-stage, DINO-compatible framework that requires no external tokenizers, enabling fair comparison that isolates the effect of the target representation from confounds of architecture and training recipe.
\textbf{(3)}~Through systematic ablation, we reveal that component importance is regime-dependent: auxiliary losses are critical from scratch but can over-constrain a pretrained backbone, yielding a practical recipe---full objective for from-scratch training, patch-only for adaptation.


% ==========================================================
% 2. RELATED WORK
% ==========================================================
\section{Related Work}
\label{sec:related}

\paragraph{Self-Distillation Methods.}
BYOL~\cite{grill2020bootstrap} demonstrated that a momentum teacher with stop-gradient enables self-supervised learning without negative pairs.
DINO~\cite{caron2021emerging} extended this with multi-crop training and centering.
DINOv2~\cite{oquab2024dinov2} scaled this recipe and introduced Sinkhorn-Knopp normalization for teacher targets along with KoLeo regularization~\cite{sablayrolles2019spreading}.
DINOv3 further refines the architecture with rotary position embeddings, register tokens, and improved training schedules; we build on DINOv3 as our backbone.
All these methods operate in continuous feature space.

\paragraph{Masked Image Modeling.}
MAE~\cite{he2022masked} reconstructs masked patches in pixel space with an asymmetric encoder-decoder.
BEiT~\cite{bao2022beit} replaces pixel targets with discrete visual tokens from a pretrained dVAE.
BEiT~v2~\cite{peng2022beitv2} refines this with a VQ-KD tokenizer but requires a two-stage pipeline.
I-JEPA~\cite{assran2023self} predicts latent representations of masked regions.
SdAE~\cite{chen2022sdae} uses EMA teacher latent features as reconstruction targets.
EVA~\cite{fang2023eva} reconstructs CLIP features at masked positions.

\paragraph{Combining Masking and Self-Distillation.}
iBOT~\cite{zhou2022ibot} combines DINO-style CLS distillation with masked patch prediction but remains in continuous space.
DINOv2~\cite{oquab2024dinov2} and its successor DINOv3 integrate DINO and iBOT losses with Sinkhorn normalization, still predicting continuous features.
MOCA~\cite{gidaris2024moca} is the most closely related prior work: it predicts soft codebook assignments at masked positions with an EMA teacher.
Despite its contribution, MOCA does not address several design choices that we show are consequential in the low-data regime: it uses softmax targets without Sinkhorn balancing, operates with only two views, does not align CLS tokens in codebook space, and was evaluated exclusively under large-scale pretraining (ImageNet).
\method{} completes this design with balanced assignments, multi-crop CLS alignment, and---crucially---targets the under-explored low-data adaptation setting where codebook structure matters most.
SIGMA~\cite{li2024sigma} independently validates Sinkhorn-guided discrete assignments for masked video modeling at ECCV~2024, reinforcing the broader value of balanced discrete targets.

\paragraph{Codebook and Clustering Methods.}
SwAV~\cite{caron2020unsupervised} introduced online clustering with Sinkhorn-Knopp for self-supervised learning on image-level features.
\method{} extends balanced assignments to both patch-level masked prediction and CLS-level cross-view learning within a single learned codebook.
We emphasize that \method{} is \emph{not} a new architecture: it adopts the standard DINO teacher-student framework and changes only the target space.
This minimal design is intentional---it enables controlled ablation that isolates the effect of discrete targets from confounds of architecture, optimizer, and data pipeline, an experimental property we consider essential for testing the codebook hypothesis.


% ==========================================================
% 3. METHOD
% ==========================================================
\section{Method}
\label{sec:method}

\method{} inherits the teacher-student multi-view protocol of DINO~\cite{caron2021emerging} but changes one critical design choice: the target space (\cref{fig:architecture}).
Instead of regressing continuous teacher features, the student predicts Sinkhorn-balanced distributions over a shared visual codebook $\mathbf{C} \in \mathbb{R}^{K \times d}$, constraining predictions to $K$ discrete categories.
The total objective combines three terms:
\begin{equation}
    \mathcal{L}_{\text{total}} = \underbrace{\mathcal{L}_{\text{mask}}}_{\text{masked codebook}} + \lambda_{\text{CLS}} \underbrace{\mathcal{L}_{\text{CLS}}}_{\text{cross-view CLS}} + \lambda_{\text{KoLeo}} \underbrace{\mathcal{L}_{\text{KoLeo}}}_{\text{uniformity}}.
    \label{eq:total_loss}
\end{equation}
$\mathcal{L}_{\text{mask}}$ (\cref{sec:kodiak_loss}) is the core objective: predict balanced codebook assignments at masked positions from teacher targets.
$\mathcal{L}_{\text{CLS}}$ (\cref{sec:cls_loss}) aligns CLS tokens across global and local crops in the same $K$-dimensional categorical space.
$\mathcal{L}_{\text{KoLeo}}$ (\cref{sec:koleo}) prevents representation collapse via a uniformity regularizer.
We first describe the shared architecture (\cref{sec:architecture}), then detail each loss.

\begin{figure}[t]
\centering
\resizebox{0.95\textwidth}{!}{\input{figures/ofr.tex}}
\caption{\textbf{\method{} architecture.} The teacher (EMA-updated) encodes unmasked global crops to produce codebook targets via Sinkhorn-Knopp balanced assignments. The student encodes masked global crops with a sparse forward pass, reconstructs the full spatial grid via a lightweight decoder, and predicts codebook assignments for masked patches. Local crops contribute CLS tokens for cross-view codebook learning.}
\label{fig:architecture}
\end{figure}

\subsection{Architecture}
\label{sec:architecture}

\paragraph{Encoders.}
Both teacher and student use a DINOv3 ViT-S/16 backbone (RoPE, register tokens, LayerScale; details in \cref{tab:architecture}).
The teacher is an EMA copy of the student and processes full, unmasked images.
The student employs MAE-style sparse forward~\cite{he2022masked}: given a binary mask $\mathbf{m} \in \{0,1\}^N$ over the $N$ spatial patches, only visible tokens enter the encoder, reducing computation proportional to the masking ratio.

\paragraph{Decoder and adapter.}
A lightweight transformer decoder (4 blocks, dim 192) reconstructs the full $N$-token spatial grid from visible encoder tokens and learnable mask tokens.
An adapter $a$ (LayerNorm + linear, $D_{\text{dec}} {\to} D$) maps the decoded features back to encoder dimensionality before the projector.

\paragraph{Projectors and visual codebook.}
MLP projectors $g_s, g_t: \mathbb{R}^D \to \mathbb{R}^d$ map patch features to codebook space and $\ell_2$-normalize them.
The codebook $\mathbf{C} \in \mathbb{R}^{K \times d}$ defines $K$ learnable visual prototypes, shared between teacher and student; logits are computed as
\begin{equation}
    \boldsymbol{\ell} = \bar{\mathbf{u}} \mathbf{C}^\top,
    \quad
    \bar{\mathbf{u}} = g(\mathbf{z}) / \norm{g(\mathbf{z})}_2,
\end{equation}
where each row of $\mathbf{C}$ represents a discrete visual concept discovered during training.

\subsection{Masked Codebook Distillation}
\label{sec:kodiak_loss}

The teacher encodes an unmasked global crop to produce patch-level logits $\boldsymbol{\ell}_t \in \mathbb{R}^{N \times K}$ over the codebook.
These logits are centered (EMA center $\mathbf{c}_{\text{patch}}$) and converted to balanced soft assignments via Sinkhorn-Knopp normalization~\cite{caron2020unsupervised}:
\begin{equation}
    \mathbf{Q} = \sinkhorn\!\left((\boldsymbol{\ell}_t - \mathbf{c}_{\text{patch}}) / \tau_t\right),
\end{equation}
where $\tau_t$ is the teacher temperature.
Sinkhorn normalization enforces an equipartition constraint across codebook entries within each batch, preventing any single prototype from dominating and ensuring diverse usage even when data is limited.

The student processes the masked view through the sparse encoder, decoder, and adapter, yielding student logits $\boldsymbol{\ell}_s \in \mathbb{R}^{N \times K}$.
The loss is a cross-entropy over masked positions $I_{\text{mask}} = \{i : m_i = 1\}$ only:
\begin{equation}
    \mathcal{L}_{\text{mask}} = -\frac{1}{|I_{\text{mask}}|}
    \sum_{i \in I_{\text{mask}}} \sum_{k=1}^{K} Q_{ik}\, \log \softmax(\boldsymbol{\ell}_{s,i})_k.
\end{equation}
By predicting categorical distributions over $K$ codebook entries---rather than regressing $d$-dimensional feature vectors---the loss imposes a discrete information bottleneck that serves as implicit regularization, combining the masked prediction paradigm of MAE~\cite{he2022masked} with the balanced assignment principle of SwAV~\cite{caron2020unsupervised}.

\subsection{Cross-View Codebook Learning}
\label{sec:cls_loss}

$\mathcal{L}_{\text{mask}}$ provides dense patch-level supervision; we complement it with an image-level objective for view-invariance.
Following DINO's multi-crop protocol~\cite{caron2021emerging}, the teacher encodes two unmasked global crops and the student processes both masked global crops and $n_\ell$ unmasked local crops, each yielding a CLS token.

A shared MLP head $h: \mathbb{R}^D \to \mathbb{R}^K$ projects CLS tokens to $K$-dimensional distributions---a separate pathway from the patch-level projector $g$ and codebook $\mathbf{C}$, but sharing the categorical dimensionality $K$ so that patch-level and image-level discrete decisions are coupled in granularity.
Teacher targets are centered (EMA center $\mathbf{c}_{\text{cls}}$) and sharpened:
$\mathbf{p}_{t,i} = \softmax\!\bigl((h(\mathbf{c}_{t,i}) - \mathbf{c}_{\text{cls}})/\tau_t\bigr)$,
while student predictions use a softer temperature $\tau_s > \tau_t$:
$\mathbf{q}_{s} = \logsoftmax\!\bigl(h(\mathbf{c}_{s})/\tau_s\bigr)$.
The loss sums over all cross-view pairs:
\begin{equation}
    \mathcal{L}_{\text{CLS}} = -\frac{1}{|P|}\sum_{(s,t) \in P} \mathbf{p}_t^\top \mathbf{q}_s,
    \label{eq:cls_loss}
\end{equation}
where $P$ enumerates global$\leftrightarrow$global and local$\to$global pairs.
As shown in \cref{sec:ablation_components}, this loss is critical from scratch but becomes optional when adapting a pretrained backbone that already exhibits view-invariance.

\subsection{KoLeo Regularization}
\label{sec:koleo}

We apply the KoLeo differential entropy estimator~\cite{sablayrolles2019spreading} as a uniformity regularizer on $\ell_2$-normalized CLS tokens $\{\bar{\mathbf{c}}_i\}_{i=1}^B$:
\begin{equation}
    \mathcal{L}_{\text{KoLeo}} =
    -\frac{1}{B} \sum_{i=1}^{B} \log \!\left( \min_{j \neq i} \norm{\bar{\mathbf{c}}_i - \bar{\mathbf{c}}_j}_2 + \epsilon \right).
\end{equation}
While Sinkhorn balancing operates at the codebook level, KoLeo acts in representation space, ensuring the encoder utilizes its full capacity under limited data diversity.
The complete training procedure is detailed in \cref{alg:kodiak} (appendix).


% ==========================================================
% 4. EXPERIMENTS
% ==========================================================
\section{Experiments}
\label{sec:experiments}

To disentangle the contributions of the training objective, the initialization, and the domain data, we evaluate \method{} under three complementary settings denoted throughout tables and figures as:
\begin{itemize}[nosep,leftmargin=*]
    \item[\basesym{}] \textbf{Foundation model baseline.}
    The off-the-shelf DINOv3 ViT-S/16 encoder pretrained on LVD-142M~\cite{oquab2024dinov2} is evaluated \emph{directly} on each target domain without any domain-specific continued pretraining.
    This reference isolates the quality of general-purpose features and provides an upper bound that domain adaptation must exceed to be considered beneficial.
    \item[\fsym{}] \textbf{From-scratch pretraining.}
    Each method is trained on the target-domain data only, starting from random initialization.
    This stress test removes the confound of pretrained features and isolates the training objective as the sole factor determining representation quality.
    \item[\csym{}] \textbf{Continued pretraining (Cont.~PT).}
    Each method initializes from the same DINOv3 foundation weights and continues self-supervised training on the target domain.
    This is the practical scenario: a practitioner adapts an existing foundation model to a specialized domain with limited unlabelled data.
\end{itemize}
Comparing \basesym{} against \csym{} reveals whether continued pretraining \emph{helps} or \emph{hurts} for a given method and domain; comparing \fsym{} against \csym{} quantifies the benefit of foundation-model initialization.

% ----------------------------------------------------------
\subsection{Experimental Setup}
\label{sec:setup}

\paragraph{Datasets.}
We evaluate on four benchmarks spanning diverse visual domains.
To keep the main paper focused, the dataset summary is deferred to the appendix (\cref{tab:datasets}).

\paragraph{Baselines.}
We compare against four self-supervised baselines, all sharing the same ViT-S/16 backbone:
DINOv3 (continuous self-distillation with iBOT-style masked prediction; the official LVD-142M checkpoint serves as the foundation model),
MAE~\cite{he2022masked} (pixel-space masked reconstruction),
I-JEPA~\cite{assran2023self} (joint-embedding predictive architecture),
and MOCA~\cite{gidaris2024moca} (discrete codebook prediction with softmax targets, no Sinkhorn balancing).

\paragraph{Evaluation protocol.}
We employ four protocols of increasing adaptation strength:
\begin{itemize}[nosep,leftmargin=*]
    \item \textbf{$k$-NN}: nearest-neighbor classification on frozen features ($k{=}20$, $\tau{=}0.07$).
    \item \textbf{Linear probing (LP)}: a linear head trained on frozen encoder features.
    \item \textbf{Low-shot classification}: 1--16 labelled examples per class, both LP and fine-tuning (\cref{fig:lowshot}).
    \item \textbf{Full fine-tuning (FT)}: end-to-end training with full label access.
\end{itemize}

\paragraph{Controlled comparison.}
All methods share the same ViT-S/16 backbone, augmentation pipeline, optimizer (AdamW, cosine schedule), batch size (128), and training budget.
For continued pretraining, all methods initialize from the same DINOv3 checkpoint.
Each baseline retains its method-specific components (\eg, MAE's asymmetric decoder, I-JEPA's context predictor) so that comparisons isolate the training objective, not the architecture.

\paragraph{Architecture scope.}
We standardize on ViT-S/16 for all comparisons.
Within an acceptable compute budget, ViT-S permits exhaustive ablation across four datasets, two training regimes, six codebook sizes, and four component variants---over 80 pretraining runs.
Smaller models are \emph{more} sensitive to the training objective than larger ones~\cite{he2022masked}, making ViT-S the strongest test of our hypothesis.
All baselines were originally developed and validated at the ViT-S scale.

% ----------------------------------------------------------
\subsection{Main Results}
\label{sec:main_results}

\begin{table}[t]
\caption{\textbf{Main results.}
$k$-NN, LP, and fine-tuning accuracy (\%) with ViT-S/16 on three target domains.
We distinguish \best{best} and \second{second-best} results within each regime.
\baseline{Gray}: \basesym\;foundation model (DINOv3 pretrained, no domain adaptation).
\fsym: from-scratch pretraining.
\csym: continued pretraining from DINOv3.}
\label{tab:contpt}
\setlength{\tabcolsep}{3pt}\renewcommand{\arraystretch}{1.05}\centering\scriptsize
\begin{tabular}{@{}l ccc ccc ccc@{}}
\toprule
& \multicolumn{3}{c}{\textbf{EuroSAT}} & \multicolumn{3}{c}{\textbf{DTD}} & \multicolumn{3}{c}{\textbf{Oxford-Pets}} \\
\cmidrule(lr){2-4} \cmidrule(lr){5-7} \cmidrule(lr){8-10}
& $k$-NN & LP & FT & $k$-NN & LP & FT & $k$-NN & LP & FT \\
\midrule
\rowcolor{baselinecolor}
\basesym\;Foundation & 94.95 & 93.06 & 97.32 & 79.86 & 79.42 & 71.12 & 92.93 & 93.23 & 86.22 \\
\midrule
\multicolumn{10}{@{}l}{\fsym\;\textit{From-scratch pretraining}} \\
DINOv3\fsym & \best{98.59} & \second{92.76} & \NA & 13.05 & \best{36.49} & 10.45 & \second{7.95} & \second{9.54} & 10.84 \\
MAE\fsym & 86.02 & 72.01 & 96.44 & 2.69 & 1.95 & 12.42 & 2.59 & 2.72 & 11.17 \\
I-JEPA\fsym & 86.35 & 82.62 & 94.97 & \second{25.16} & \second{23.04} & 24.45 & 7.30 & 8.67 & 9.37 \\
MOCA\fsym & 94.63 & 93.86 & \best{98.37} & 25.13 & 26.06 & \second{34.93} & 6.36 & 10.78 & \second{15.65} \\
\rowours
\oursshort\fsym & \second{96.94} & \best{94.56} & \second{97.57} & \best{46.03} & \NA & \best{49.69} & \best{10.12} & \best{12.60} & \best{21.80} \\
\midrule
\multicolumn{10}{@{}l}{\csym\;\textit{Continued pretraining}} \\
DINOv3\csym & \second{97.52} & 94.33 & 98.31 & 60.41 & 42.84 & 62.44 & 67.90 & 20.64 & 80.26 \\
MAE\csym & 80.61 & 65.25 & 97.74 & 29.14 & 36.11 & 73.39 & 21.27 & 32.93 & 88.16 \\
I-JEPA\csym & 96.26 & 96.45 & 98.69 & \best{77.68} & \second{76.84} & 75.99 & 78.17 & 88.24 & 88.12 \\
MOCA\csym & 97.15 & \second{97.28} & \second{98.77} & \second{77.00} & \best{76.88} & \second{77.03} & \best{83.75} & \second{90.58} & \second{89.58} \\
\rowours
\ours\csym & \best{98.59} & \best{97.89} & \best{98.93} & 69.59 & 75.20 & \best{78.00} & \second{80.00} & \best{91.03} & \best{91.61} \\
\bottomrule
\end{tabular}
\end{table}

\cref{tab:contpt} reports results across $k$-NN, LP, and fine-tuning for all three evaluation regimes.
We highlight the key observations:
\begin{itemize}[nosep,leftmargin=*]
    \item \textbf{Continuous objectives degrade under low data.}
    Continuing DINOv3 pretraining degrades \emph{all three} protocols on DTD and Oxford-Pets.
    The collapse is most severe for frozen features: LP on Oxford-Pets drops by $>$72\,pp (from 93.2\% to 20.6\%), confirming that continuous self-distillation destroys representational structure.
    MAE follows a similar pattern ($>$$-$27\,pp LP on every dataset), though it partially recovers under fine-tuning---consistent with its known tendency to learn features that require end-to-end adaptation.

    \item \textbf{Discrete targets preserve and improve the backbone.}
    Both discrete-codebook methods (\method{} and MOCA) improve over the foundation model on all three datasets under fine-tuning, with \method{} achieving gains of $+$6.9\,pp (DTD) and $+$5.4\,pp (Oxford-Pets).
    Frozen-feature quality is largely preserved: \method{} retains 91.0\% LP on Oxford-Pets vs.\ 20.6\% for DINOv3\csym{}.
    I-JEPA also preserves features well (joint-embedding design), but \method{} achieves the highest FT accuracy on two of three datasets.

    \item \textbf{From-scratch stress test.}
    Training from random initialization (upper block) isolates the objective from initialization quality.
    \method{} outperforms the next-best method by $>$$+$25\,pp on DTD fine-tuning (49.7\% vs.\ I-JEPA's 24.5\%), while DINOv3 collapses to 10.5\%.
    On frozen features, \method{} reaches 46.0\% $k$-NN on DTD (nearly $2{\times}$ I-JEPA) and leads LP on EuroSAT (94.6\%) and Oxford-Pets (12.6\%).
    DINOv3 and MAE fail to learn linearly separable features on DTD and Oxford-Pets, with LP near random chance.
\end{itemize}

\paragraph{Low-shot classification (\cref{fig:lowshot}).}
We evaluate all continued-pretraining methods with 1--16 labelled examples per class under both fine-tuning and linear probing.
The \globe\;Foundation baseline (off-the-shelf DINOv3, no domain adaptation) isolates the contribution of continued pretraining from label-efficiency of the backbone itself.

\paragraph{Fine-tuning~(a).}
\method{} achieves the highest accuracy at every shot budget on all three datasets.
The advantage is most pronounced at one shot per class: 59.0\% on EuroSAT (nearly $+$25\,pp over the foundation baseline), 20.7\% on DTD ($+$8\,pp over I-JEPA).
At 16 shots the gap narrows but remains significant: $+$5.7\,pp on EuroSAT, $+$4.7\,pp on DTD, and $+$7.2\,pp on Oxford-Pets over the second-best method.

\paragraph{Linear probing~(b).}
\method{} leads on EuroSAT across all shot budgets (up to $+$31.9\,pp over \globe\;Foundation at 1~shot) and on Oxford-Pets (up to $+$16.6\,pp at 1~shot).
The foundation baseline is competitive on DTD and Oxford-Pets, occasionally matching adapted methods---suggesting that DINOv3's general-purpose features already capture texture and fine-grained patterns useful for these tasks.
At higher shot budgets MOCA overtakes on EuroSAT from 4~shots onward.
MAE and the domain-adapted DINOv3 trail by a wide margin under both protocols, confirming that continuous-target representations are less amenable to rapid adaptation with limited supervision.

\begin{figure}[t]
\centering
\includegraphics[width=\linewidth]{figures/kshot_classification.pdf}
\caption{\textbf{Low-shot classification} (continued pretraining, \csym).
Top-1 accuracy vs.\ shots per class (1--16) on three domains.
(a)~End-to-end fine-tuning.
(b)~Linear probing (frozen backbone).
Mean$\pm$std over 3 label seeds.
\basesym\;\emph{Foundation}: off-the-shelf DINOv3 encoder, no domain adaptation.
\method{} leads at nearly all shot budgets.}
\label{fig:lowshot}
\end{figure}

% ----------------------------------------------------------
\subsection{Representation Drift Analysis}
\label{sec:cka}

We quantify how much continued pretraining changes the foundation representation using linear CKA~\cite{kornblith2019similarity} computed layer-by-layer between the off-the-shelf DINOv3 encoder and each adapted encoder, using the same 1\,024 sampled test images for all methods on each dataset.
For each of the 12 transformer blocks, we extract the post-block CLS token embedding (index~0 of the output sequence, excluding register and patch tokens) and compute pairwise linear CKA between the foundation and each adapted encoder.
We define a scalar \emph{drift index} $D = 1 - \overline{\text{CKA}}_{\text{diag}}$, where the bar denotes the mean over the 12 transformer blocks.

\begin{figure}[t]
\centering
\begin{subfigure}[b]{0.48\textwidth}
    \centering
    \includegraphics[width=\textwidth]{figures/cka_layerwise_main.pdf}
    \caption{Per-block CKA curves (DTD and Oxford-Pets).}
    \label{fig:cka_curves}
\end{subfigure}
\hfill
\begin{subfigure}[b]{0.48\textwidth}
    \centering
    \includegraphics[width=\textwidth]{figures/cka_heatmaps_main.pdf}
    \caption{Full layer$\times$layer CKA heatmaps, all methods.}
    \label{fig:cka_heatmaps}
\end{subfigure}
\caption{\textbf{Representation drift under continued pretraining.}
(a)~Layer-wise CKA between the foundation and each adapted model on DTD and Oxford-Pets.
DINOv3 and MAE drift heavily in mid-to-late layers; \method{} and I-JEPA preserve the backbone.
(b)~Full layer$\times$layer CKA heatmaps.
Drift index $D$ (inset) summarizes overall change; lower is better.}
\label{fig:cka}
\end{figure}

\cref{fig:cka} reveals that on DTD and Oxford-Pets, the two datasets where continued DINOv3 pretraining \emph{hurts}, DINOv3 and MAE undergo severe representational drift ($D > 0.48$), with per-block CKA dropping below 0.2 in late layers.
\method{} preserves high layer-wise alignment ($D \approx 0.11$--$0.14$), and I-JEPA behaves similarly ($D \approx 0.10$--$0.15$), consistent with its joint-embedding design.
Methods with lower CKA preservation consistently suffer larger LP drops, confirming that representational drift, not suboptimal hyperparameters, explains the collapse in \cref{tab:contpt}.

% ----------------------------------------------------------
\subsection{Ablation Studies}
\label{sec:ablations}
\label{sec:ablation_components}
\label{sec:codebook_size}

\begin{table}[t]
\caption{\textbf{Design analysis.}
(a)~Component ablation: FT accuracy (\%) under both regimes; \textcolor{kodiakgreen}{green}/\textcolor{kodiakred}{red} deltas vs.\ full \oursshort{}.
(b)~Codebook size ($K$) sweep.
\baseline{Gray}: default ($K{=}128$).
\fsym: from scratch. \csym: continued PT.}
\label{tab:design}
\begin{subtable}[t]{0.48\textwidth}
\caption{Component ablation}
\label{tab:ablations}
\tablestyle{2pt}{1.05}
\begin{tabular}{@{}lccc@{}}
\toprule
Variant & EuroSAT & DTD & Ox-Pets \\
\midrule
\multicolumn{4}{@{}l}{\fsym\;\textit{From scratch}} \\
\baseline{Full \oursshort} & \baseline{\best{97.57}} & \baseline{\best{49.69}} & \baseline{21.80} \\
w/o Sinkhorn & 97.57 & 48.24 \drop{1.5} & \best{24.09} \gain{2.3} \\
w/o $\mathcal{L}_{\text{CLS}}$ & 95.99 \drop{1.6} & 25.36 \drop{24.3} & 19.52 \drop{2.3} \\
w/o $\mathcal{L}_{\text{KoLeo}}$ & 96.83 \drop{0.7} & 25.01 \drop{24.7} & 14.32 \drop{7.5} \\
\midrule
\multicolumn{4}{@{}l}{\csym\;\textit{Continued PT}} \\
\baseline{Full \oursshort} & \baseline{\best{98.93}} & \baseline{78.00} & \baseline{91.61} \\
w/o Sinkhorn & 98.60 \drop{0.3} & 77.50 \drop{0.5} & 91.60 \\
w/o $\mathcal{L}_{\text{CLS}}$ & 98.38 \drop{0.6} & \best{78.80} \gain{0.8} & \best{93.13} \gain{1.5} \\
w/o $\mathcal{L}_{\text{KoLeo}}$ & 98.68 \drop{0.3} & 78.68 \gain{0.7} & 91.21 \drop{0.4} \\
\bottomrule
\end{tabular}
\end{subtable}
\hfill
\begin{subtable}[t]{0.48\textwidth}
\caption{Codebook size ($K$)}
\label{tab:codebook_size}
\tablestyle{2pt}{1.05}
\begin{tabular}{@{}rccc@{}}
\toprule
$K$ & EuroSAT & DTD & Ox-Pets \\
\midrule
\multicolumn{4}{@{}l}{\fsym\;\textit{From scratch}} \\
64   & 97.32 & 47.65 & 21.02 \\
\rowcolor{baselinecolor}
128  & 97.57 & \best{49.69} & 21.80 \\
256  & 97.29 & 48.72 & 21.77 \\
512  & 96.89 & 46.68 & \best{24.61} \\
1024 & 97.14 & 47.15 & 24.51 \\
4096 & \best{97.73} & 47.27 & 21.46 \\
\midrule
\multicolumn{4}{@{}l}{\csym\;\textit{Continued PT}} \\
64   & 98.67 & 77.95 & 91.12 \\
\rowcolor{baselinecolor}
128  & \best{98.93} & 78.00 & 91.61 \\
256  & 98.61 & 77.65 & 91.73 \\
512  & 98.56 & \best{78.65} & 92.19 \\
1024 & 98.60 & 78.15 & \best{92.24} \\
4096 & 98.41 & 78.62 & 92.10 \\
\bottomrule
\end{tabular}
\end{subtable}
\end{table}

\paragraph{Component importance is regime-dependent (\cref{tab:ablations}).}
From scratch, removing either $\mathcal{L}_{\text{CLS}}$ or $\mathcal{L}_{\text{KoLeo}}$ collapses DTD by more than $-$24\,pp, confirming that auxiliary losses provide essential regularization without pretrained features.
Under continued pretraining, the picture reverses: the codebook objective alone carries most of the gain, and dropping $\mathcal{L}_{\text{CLS}}$ even \emph{improves} Oxford-Pets by $+$1.5\,pp.
This suggests a practical recipe: full objective for from-scratch training, patch-only for adaptation.

\paragraph{Codebook size (\cref{tab:codebook_size}).}
Performance is robust across $K \in [64, 4096]$: under continued pretraining, the spread is less than 1\,pp on all datasets.
We use $K{=}128$ as default.


% ==========================================================
% 5. DISCUSSION
% ==========================================================
\section{Discussion}
\label{sec:discussion}

\paragraph{Why discrete targets enable safe adaptation.}
Our CKA analysis (\cref{sec:cka}) provides mechanistic insight: continuous self-distillation causes severe representational drift ($D > 0.48$ on DTD and Oxford-Pets), progressively destroying the backbone's linear separability in mid-to-late transformer layers.
Discrete codebook targets prevent this drift ($D \approx 0.11$--$0.14$) while retaining enough flexibility to absorb domain-specific structure.
The key mechanism is that forcing predictions through a categorical bottleneck of $K$ entries limits the degrees of freedom available for overfitting, effectively regularizing adaptation without explicit constraints on the feature space.
This explains why \method{} improves the foundation model on datasets where the standard DINO objective degrades it (\cref{tab:contpt}).

\paragraph{The role of Sinkhorn balancing.}
The discrete bottleneck alone is necessary but not sufficient: without balanced assignments, the codebook degenerates to a handful of active entries.
Our utilization analysis (\cref{tab:code_util}) shows that removing Sinkhorn leads to severe collapse (\eg, only 11 of 128 codes active on EuroSAT from scratch), even when fine-tuning accuracy appears unaffected.
This silent failure mode is dangerous in practice: the codebook loses its capacity to represent diverse visual concepts, undermining the very information bottleneck that makes discrete targets effective.

\paragraph{Regime-dependent design.}
The most practically relevant finding from our ablations (\cref{tab:ablations}) is that the optimal recipe depends on initialization.
From scratch, the auxiliary losses ($\mathcal{L}_{\text{CLS}}$, $\mathcal{L}_{\text{KoLeo}}$) are essential---removing either collapses performance by over $-$24\,pp on DTD---because the encoder lacks the geometric priors needed to learn from the codebook objective alone.
Under continued pretraining, these losses become redundant or mildly harmful, since the pretrained backbone already provides view-invariance and representational diversity.
This yields a simple practitioner guideline: \emph{full objective for from-scratch training; codebook-only for foundation-model adaptation.}

\paragraph{Scope and future directions.}
We deliberately restrict our study to ViT-S/16, where objective choice most visibly separates methods and the compute budget permits exhaustive ablation.
We expect the gap between discrete and continuous targets to narrow at larger model scales, where increased capacity can compensate for suboptimal objectives~\cite{he2022masked}.
A complementary open question is the data-efficiency boundary: our results show clear gains below 27K images (EuroSAT) but the threshold above which discrete targets lose their advantage remains to be characterized.
Finally, the strong performance of I-JEPA under adaptation (\cref{sec:cka}) suggests that joint-embedding objectives may share structural properties with discrete targets; disentangling the roles of the categorical bottleneck, balanced assignments, and the absence of pixel-level regression could inform the design of next-generation SSL objectives.


% ==========================================================
% 6. CONCLUSION
% ==========================================================
\section{Conclusion}
\label{sec:conclusion}

Continuing standard SSL objectives on small domain datasets can \emph{degrade} the very representations they aim to improve.
Our controlled study confirms that this failure tracks representational drift in the continuous target space, not suboptimal hyperparameters.
Replacing continuous feature regression with discrete, Sinkhorn-balanced codebook assignments transforms adaptation from a fragile operation into a reliable one.

\method{} instantiates this principle in a single-stage, DINO-compatible framework requiring no external tokenizers.
Under continued pretraining, it improves the foundation model on every dataset tested, while the standard DINO objective degrades two out of three.
From scratch, it substantially outperforms all baselines, confirming that discrete targets provide a stronger learning signal independent of initialization.
Crucially, the optimal recipe is regime-dependent---full objective from scratch, codebook-only for adaptation---offering practitioners a simple guideline for their data budget.
We hope the ``safe adaptation'' perspective motivates further investigation into the structural properties of SSL target spaces.


% ==========================================================
% APPENDIX
% ==========================================================
\appendix
\section*{Appendix}
\addcontentsline{toc}{section}{Appendix}

% ----------------------------------------------------------
\section{Training Procedure}
\label{sec:app_algorithm}

\cref{alg:kodiak} details the complete \method{} training loop for one minibatch, including the teacher forward pass (unmasked global crops), the student sparse forward (masked global crops), local crop encoding (CLS tokens only), and loss computation.

\begin{algorithm}[ht]
\caption{\method{} Pretraining}
\label{alg:kodiak}
\small
\begin{algorithmic}[1]
\Require Global crops $\{x^g_1, x^g_2\}$; local crops $\{x^\ell_j\}_{j=1}^{n_\ell}$; masks $\{m_1, m_2\}$
\Require Student $f_s, d, a, g_s$; Teacher $f_t, g_t$ (EMA); Shared: $h$, $\mathbf{C} \in \mathbb{R}^{K \times d}$
\For{each minibatch}
    \State \textcolor{kodiakblue}{\textbf{// Teacher: full forward on unmasked global crops}}
    \For{$i \in \{1,2\}$}
        \State $(\mathbf{c}_{t,i}, \mathbf{z}_{t,i}) \gets f_t(x^g_i)$
        \State $\boldsymbol{\ell}^{\text{patch}}_{t,i} \gets \mathrm{norm}(g_t(\mathbf{z}_{t,i})) \cdot \mathbf{C}^\top$
    \EndFor
    \State \textcolor{kodiakblue}{\textbf{// Student: sparse forward on masked global crops}}
    \For{$i \in \{1,2\}$}
        \State $(\mathbf{c}_{s,i}, \mathbf{z}^{\text{vis}}_{s,i}) \gets f_s(x^g_i, m_i)$; \quad $\hat{\mathbf{z}}_{s,i} \gets d(\mathbf{z}^{\text{vis}}_{s,i}, m_i)$
        \State $\boldsymbol{\ell}^{\text{patch}}_{s,i} \gets \mathrm{norm}(g_s(a(\hat{\mathbf{z}}_{s,i}))) \cdot \mathbf{C}^\top$
    \EndFor
    \State \textcolor{kodiakblue}{\textbf{// Student: local crops (CLS only)}}
    \For{$j = 1, \ldots, n_\ell$}
        \State $(\mathbf{c}^\ell_{s,j}, \_) \gets f_s(x^\ell_j)$
    \EndFor
    \State \textcolor{kodiakgreen}{\textbf{// Losses}}
    \State $\mathbf{Q}_i \gets \sinkhorn((\boldsymbol{\ell}^{\text{patch}}_{t,i} - \mathbf{c}_{\text{patch}})/\tau_t)$ \quad for $i \in \{1,2\}$
    \State $\mathcal{L}_{\text{mask}} \gets \tfrac{1}{2}\sum_{i} \mathrm{CE}(\mathbf{Q}_i,\, \logsoftmax(\boldsymbol{\ell}^{\text{patch}}_{s,i}))$ \quad (masked patches only)
    \State $\mathcal{L}_{\text{CLS}} \gets$ cross-view CE via head $h$ (\cref{eq:cls_loss})
    \State $\mathcal{L}_{\text{KoLeo}} \gets$ KoLeo on all student CLS tokens
    \State $\mathcal{L} \gets \mathcal{L}_{\text{mask}} + \lambda_{\text{CLS}}\,\mathcal{L}_{\text{CLS}} + \lambda_{\text{KoLeo}}\,\mathcal{L}_{\text{KoLeo}}$
    \State Update student (gradient descent); update teacher (EMA)
\EndFor
\end{algorithmic}
\end{algorithm}

% ----------------------------------------------------------
\section{Dataset Summary}
\label{sec:app_datasets}

\begin{table}[ht]
\caption{\textbf{Evaluation datasets.}
Train/val/test splits follow official partitions where available.
$^{*}$~Val carved from train (10\%).}
\label{tab:datasets}
\tablestyle{2pt}{1.05}
\begin{tabular}{@{}llrr rrr@{}}
\toprule
Dataset & Domain & Classes & Images & Train & Val & Test \\
\midrule
EuroSAT~\cite{helber2019eurosat} & Remote sensing & 10 & 27{,}000 & 16{,}200 & 5{,}400 & 5{,}400 \\
DTD~\cite{cimpoi2014dtd} & Texture & 47 & 5{,}640 & 1{,}880 & 1{,}880 & 1{,}880 \\
Oxford-Pets~\cite{parkhi2012oxfordpets} & Fine-grained & 37 & 7{,}349 & 3{,}680$^{*}$ & & 3{,}669 \\
CIFAR-100~\cite{krizhevsky2009cifar} & Object recognition & 100 & 60{,}000 & 50{,}000$^{*}$ & & 10{,}000 \\
\bottomrule
\end{tabular}
\end{table}

\paragraph{EuroSAT~\cite{helber2019eurosat}.}
EuroSAT contains 27{,}000 geo-referenced Sentinel-2 satellite image patches at 10\,m ground sampling distance, each $64 \times 64$ pixels covering 10 land use and land cover classes (\eg, industrial, residential, forest, annual crop).
The dataset captures scenes across 34 European countries, providing substantial geographic diversity despite its moderate size.
We use the community-standard 60/20/20 train/val/test partition from TorchGeo, which stratifies by class to ensure balanced evaluation.
EuroSAT represents the ``medium-sized'' end of our benchmark: large enough to train reasonable representations from scratch, yet small enough that continuous SSL objectives already benefit from domain adaptation.

\paragraph{DTD~\cite{cimpoi2014dtd}.}
The Describable Textures Dataset (DTD) comprises 5{,}640 texture images collected from the wild, organized into 47 perceptual categories (\eg, braided, bubbly, woven, zigzagged) with 120 images per class.
Images vary greatly in scale, viewpoint, and illumination, making texture recognition challenging even for models pretrained on object-centric datasets.
The official release provides 10 predefined folds, each splitting the data into equal thirds (1{,}880 train / 1{,}880 val / 1{,}880 test); we report results on the first fold, following recent practice~\cite{grill2020bootstrap}.
At only 5{,}640 images, DTD is the smallest dataset in our benchmark and the one where continuous pretraining objectives fail most severely.

\paragraph{Oxford-Pets~\cite{parkhi2012oxfordpets}.}
The Oxford-IIIT Pet Dataset consists of 7{,}349 images spanning 37 breeds of cats and dogs, with approximately 200 images per class.
Intra-class variation is high---images range from close-up portraits to full-body shots in cluttered backgrounds---while inter-class differences can be subtle (distinguishing between visually similar breeds).
The official split provides 3{,}680 trainval and 3{,}669 test images; we carve a 10\% validation set from the training portion.
Oxford-Pets tests whether adapted representations retain the fine-grained discriminative structure needed for breed-level recognition.

\paragraph{CIFAR-100~\cite{krizhevsky2009cifar}.}
CIFAR-100 contains 60{,}000 colour images of size $32 \times 32$ pixels distributed evenly across 100 object categories (500 training and 100 test images per class), grouped into 20 superclasses.
Although small in resolution, its broad category coverage and balanced class distribution make it a widely used benchmark for evaluating representation quality.
The official split provides 50{,}000 training and 10{,}000 test images; we hold out 10\% of the training set for validation.
CIFAR-100 serves as a control point in our study: at 60{,}000 images it is an order of magnitude larger than DTD or Oxford-Pets, probing whether the advantages of discrete targets persist when data is less scarce.

% ----------------------------------------------------------
\section{Model and Training Configuration}
\label{sec:app_config}

\begin{table}[ht]
\caption{\textbf{Model architecture.}}
\label{tab:architecture}
\tablestyle{2pt}{1.05}
\begin{tabular}{@{}ll@{}}
\toprule
\multicolumn{2}{@{}l}{\textit{Encoder (ViT-S/16, DINOv3)}} \\[2pt]
Embedding dim $D$ & 384 \\
Depth $L$ / Heads $H$ & 12 / 6 \\
Patch size / Image size & $16{\times}16$ / $256{\times}256$ \\
MLP ratio & 4.0 \\
Position embedding & RoPE ($\theta{=}100$) \\
Register tokens $R$ & 4 \\
Drop path rate & 0.1 \\
LayerScale init & $10^{-5}$ \\
\midrule
\multicolumn{2}{@{}l}{\textit{Decoder (student only)}} \\[2pt]
Embedding dim $D_{\text{dec}}$ & 192 \\
Depth $L_{\text{dec}}$ / Heads $H_{\text{dec}}$ & 4 / 6 \\
Position embedding & learnable (bilinear interpolation) \\
Adapter $a$ & LayerNorm + Linear($D_{\text{dec}} {\to} D$) \\
\midrule
\multicolumn{2}{@{}l}{\textit{Patch projectors $g_s, g_t$ (3-layer MLP)}} \\[2pt]
Hidden dim / Output dim $d$ & 2048 / 256 \\
\midrule
\multicolumn{2}{@{}l}{\textit{CLS head $h$ (3-layer MLP)}} \\[2pt]
Architecture & $D {\to} 2048 {\to} 2048 {\to} K$ \\
\midrule
\multicolumn{2}{@{}l}{\textit{Codebook}} \\[2pt]
Entries $K$ / Dimension & 128 / 256 \\
\bottomrule
\end{tabular}
\end{table}

\begin{table}[ht]
\caption{\textbf{Training configuration.} Regime-specific values shown in separate columns; shared values span both.}
\label{tab:config}
\tablestyle{2pt}{1.05}
\begin{tabular}{@{}l cc@{}}
\toprule
& \textbf{From scratch} & \textbf{Cont.~PT} \\
\midrule
\multicolumn{3}{@{}l}{\textit{Pretraining}} \\[2pt]
Epochs & 500 & 100 \\
Batch size & \multicolumn{2}{c}{128} \\
Learning rate & $10^{-4}$ & $10^{-5}$ \\
LR schedule & \multicolumn{2}{c}{cosine} \\
LR warmup & 10 epochs & 5 epochs \\
Weight decay & \multicolumn{2}{c}{$0.04 \to 0.4$ (cosine)} \\
Masking ratio & \multicolumn{2}{c}{$[0.1, 0.5]$ (block masking)} \\
Multi-crop & \multicolumn{2}{c}{2 global ($256^2$) + 8 local ($112^2$)} \\
$\lambda_{\text{CLS}} / \lambda_{\text{KoLeo}}$ & \multicolumn{2}{c}{1.0 / 0.1} \\
Student temp. $\tau_s$ & \multicolumn{2}{c}{0.1} \\
Teacher temp. $\tau_t$ & \multicolumn{2}{c}{$0.04 \to 0.07$} \\
Temp. warmup & 30 epochs & 10 epochs \\
Center momentum & \multicolumn{2}{c}{0.9} \\
Sinkhorn iterations & \multicolumn{2}{c}{3} \\
EMA momentum & $0.992 \to 1.0$ & $0.996 \to 1.0$ \\
Freeze last layer & 1 epoch & 3 epochs \\
\midrule
\multicolumn{3}{@{}l}{\textit{Fine-tuning}} \\[2pt]
Epochs & 100 & 50 \\
Batch size & \multicolumn{2}{c}{128} \\
Learning rate & \multicolumn{2}{c}{$10^{-4}$} \\
Backbone LR scale & \multicolumn{2}{c}{0.1} \\
Early stopping & \multicolumn{2}{c}{patience 10} \\
Optimizer & \multicolumn{2}{c}{AdamW} \\
Precision & \multicolumn{2}{c}{bf16-mixed} \\
\midrule
\multicolumn{3}{@{}l}{\textit{Linear evaluation}} \\[2pt]
Epochs & \multicolumn{2}{c}{50} \\
Learning rate & \multicolumn{2}{c}{$10^{-3}$} \\
Backbone & \multicolumn{2}{c}{frozen} \\
\midrule
\multicolumn{3}{@{}l}{\textit{$k$-NN evaluation}} \\[2pt]
$k$ / temperature & \multicolumn{2}{c}{20 / 0.07} \\
\bottomrule
\end{tabular}
\end{table}

% ----------------------------------------------------------
\section{Codebook Utilization Tables}
\label{sec:app_codeutil}

\begin{table}[ht]
\caption{\textbf{Codebook utilization: component ablations} ($K{=}128$).
Active: unique entries assigned per epoch; Entropy: CLS assignment entropy (nats).}
\label{tab:code_util}
\tablestyle{2pt}{1.05}
\begin{tabular}{@{}l cc cc cc cc@{}}
\toprule
& \multicolumn{2}{c}{\cellcolor{oursrow} Full \oursshort{}} & \multicolumn{2}{c}{w/o Sinkhorn} & \multicolumn{2}{c}{w/o $\mathcal{L}_{\text{CLS}}$} & \multicolumn{2}{c}{w/o $\mathcal{L}_{\text{KoLeo}}$} \\
\cmidrule(lr){2-3} \cmidrule(lr){4-5} \cmidrule(lr){6-7} \cmidrule(lr){8-9}
Dataset & Active & Entr. & Active & Entr. & Active & Entr. & Active & Entr. \\
\midrule
\multicolumn{9}{@{}l}{\fsym{}\;\textit{From-scratch pretraining}} \\[2pt]
EuroSAT & \cellcolor{oursrow!50} 49 & \cellcolor{oursrow!50} 0.48 & \textbf{11} & 0.48 & 77 & -- & 93 & 0.73 \\
DTD & \cellcolor{oursrow!50} 91 & \cellcolor{oursrow!50} 0.34 & 78 & 0.33 & 85 & -- & 99 & 3.85 \\
Ox-Pets & \cellcolor{oursrow!50} 64 & \cellcolor{oursrow!50} 3.38 & 93 & 3.08 & 75 & -- & 70 & 4.85 \\
\midrule
\multicolumn{9}{@{}l}{\csym{}\;\textit{Continued pretraining (from DINOv3)}} \\[2pt]
EuroSAT & \cellcolor{oursrow!50} 90 & \cellcolor{oursrow!50} 1.26 & \textbf{18} & 1.37 & 105 & -- & 32 & 1.46 \\
DTD & \cellcolor{oursrow!50} 85 & \cellcolor{oursrow!50} 2.94 & 38 & 2.91 & 102 & -- & 59 & 2.44 \\
Ox-Pets & \cellcolor{oursrow!50} 105 & \cellcolor{oursrow!50} 3.35 & 91 & 3.34 & 110 & -- & 82 & 3.22 \\
\bottomrule
\end{tabular}
\end{table}

\begin{table}[ht]
\caption{\textbf{Codebook utilization: size sweep.}
Active codes and CLS entropy (nats) at final epoch for varying $K$; max entropy $= \ln K$.}
\label{tab:code_util_size}
\tablestyle{2pt}{1.05}
\begin{tabular}{@{}r cc cc cc@{}}
\toprule
& \multicolumn{2}{c}{EuroSAT} & \multicolumn{2}{c}{DTD} & \multicolumn{2}{c}{Oxford-Pets} \\
\cmidrule(lr){2-3} \cmidrule(lr){4-5} \cmidrule(lr){6-7}
$K$ & Active & Entr. & Active & Entr. & Active & Entr. \\
\midrule
\multicolumn{7}{@{}l}{\fsym{}\;\textit{From-scratch pretraining}} \\[2pt]
64   & 53  & 0.22 & 62  & 0.34 & 60  & 3.14 \\
128  & 49  & 0.48 & 115 & 0.36 & 64  & 3.39 \\
256  & 86  & 0.33 & 178 & 0.34 & 174 & 3.71 \\
512  & 210 & 0.49 & 156 & 0.37 & 187 & 2.70 \\
1024 & 150 & 0.35 & 227 & 0.37 & 468 & 2.43 \\
2048 & 337 & 0.38 & 175 & 0.40 & 262 & 2.96 \\
4096 & 194 & 0.33 & 268 & 0.44 & 256 & 8.32 \\
\midrule
\multicolumn{7}{@{}l}{\csym{}\;\textit{Continued pretraining (from DINOv3)}} \\[2pt]
64   & 56  & 1.00 & 49  & 1.67 & 58  & 2.48 \\
128  & 89  & 1.27 & 84  & 2.94 & 105 & 3.35 \\
256  & 118 & 1.60 & 138 & 5.16 & 178 & 5.10 \\
512  & 149 & 2.89 & 188 & 5.97 & 294 & 6.04 \\
1024 & 238 & 3.64 & 256 & 6.76 & 481 & 6.77 \\
2048 & 389 & 6.00 & 359 & 7.10 & 628 & 7.48 \\
4096 & 614 & 8.19 & 423 & 8.17 & 796 & 8.17 \\
\bottomrule
\end{tabular}
\end{table}

% ----------------------------------------------------------
\section{Training Dynamics and Data Efficiency}
\label{sec:app_convergence}

Beyond final accuracy, training dynamics reveal how quickly and stably each method learns under limited data.
\cref{fig:convergence_dtd} reports linear probing validation accuracy on DTD as a function of the training epoch under both regimes.
Curves are truncated at peak accuracy (star marker); shaded regions show $\pm$1 std over 3 seeds.
All methods share the same ViT-S/16 backbone, optimizer (AdamW), augmentation pipeline, batch size (128), and training budget per regime, ensuring differences reflect only the training objective.

\paragraph{Faster improvement early in training.}
Under continued pretraining~(b), \method{} reaches strong validation accuracy within the first 10--15 epochs, while DINOv3 and MAE lag far behind across the entire schedule.
This rapid improvement reflects the data efficiency of discrete codebook targets: fewer passes over the small dataset are needed to learn useful frozen features.

\paragraph{Stable plateau without oscillation.}
\method{}'s curve saturates smoothly and maintains its peak without the oscillations visible in continuous-target baselines.
This supports the claim that discrete, balanced assignments reduce representational drift and yield more robust optimization under limited data diversity.

\paragraph{Final performance is not the only story.}
Under continued pretraining, MOCA achieves a marginally higher peak LP accuracy than \method{}, but \method{} attains near-peak performance substantially earlier.
From scratch~(a), \method{} converges to a peak well above all baselines, while DINOv3 and MAE fail to learn linearly separable features at all, consistent with their near-random LP scores in \cref{tab:contpt}.
These dynamics reinforce our central theme: \method{} provides \emph{safe, data-efficient adaptation}, not just better final numbers.

\begin{figure}[ht]
\centering
\includegraphics[width=\linewidth]{figures/convergence_dtd_lp.pdf}
\caption{\textbf{Training dynamics under low-data adaptation.}
Linear probing validation accuracy (\%) vs.\ training epoch on DTD under (a)~from-scratch and (b)~continued pretraining.
Curves report mean over 3 seeds; shaded region shows standard deviation.
All methods share the same ViT-S/16 backbone, optimizer, augmentations, batch size, and training budget.
\method{} improves rapidly and plateaus smoothly, while continuous-target objectives (DINOv3, MAE) learn slowly and converge to substantially lower accuracy.
MOCA reaches a slightly higher final LP accuracy in~(b), but \method{} attains near-peak performance earlier.}
\label{fig:convergence_dtd}
\end{figure}


% ==========================================================
% ACKNOWLEDGEMENTS
% ==========================================================
\section*{Acknowledgements}
% TODO FINAL: Insert acknowledgements here.


% ==========================================================
% REFERENCES
% ==========================================================
%
% BibTeX users should specify bibliography style 'splncs04'.
% References will then be sorted and formatted in the correct style.
%
\bibliographystyle{splncs04}
% \bibliography{main}
%
% If using inline references instead of a .bib file:
\begin{thebibliography}{99}

\bibitem{assran2023self}
Assran, M., et al.:
Self-supervised learning from images with a joint-embedding predictive architecture.
In: CVPR (2023)

\bibitem{bao2022beit}
Bao, H., Dong, L., Wei, F.:
BEiT: BERT pre-training of image transformers.
In: ICLR (2022)

\bibitem{caron2020unsupervised}
Caron, M., et al.:
Unsupervised learning of visual features by contrasting cluster assignments.
In: NeurIPS (2020)

\bibitem{caron2021emerging}
Caron, M., et al.:
Emerging properties in self-supervised vision transformers.
In: ICCV (2021)

\bibitem{chen2022sdae}
Chen, Y., et al.:
SdAE: Self-distillated masked autoencoder.
In: ECCV (2022)

\bibitem{cimpoi2014dtd}
Cimpoi, M., Maji, S., Kokkinos, I., Mohamed, S., Vedaldi, A.:
Describing textures in the wild.
In: CVPR (2014)

\bibitem{cuturi2013sinkhorn}
Cuturi, M.:
Sinkhorn distances: Lightspeed computation of optimal transport.
In: NeurIPS (2013)

\bibitem{fang2023eva}
Fang, Y., et al.:
EVA: Exploring the limits of masked visual representation learning at scale.
In: CVPR (2023)

\bibitem{gidaris2024moca}
Gidaris, S., et al.:
MOCA: Self-supervised representation learning by predicting masked online codebook assignments.
TMLR (2024)

\bibitem{grill2020bootstrap}
Grill, J.-B., et al.:
Bootstrap your own latent: A new approach to self-supervised learning.
In: NeurIPS (2020)

\bibitem{he2022masked}
He, K., et al.:
Masked autoencoders are scalable vision learners.
In: CVPR (2022)

\bibitem{helber2019eurosat}
Helber, P., Bischke, B., Dengel, A., Borth, D.:
EuroSAT: A novel dataset and deep learning benchmark for land use and land cover classification.
IEEE J-STARS (2019)

\bibitem{kornblith2019similarity}
Kornblith, S., Norouzi, M., Lee, H., Hinton, G.:
Similarity of neural network representations revisited.
In: ICML (2019)

\bibitem{krizhevsky2009cifar}
Krizhevsky, A.:
Learning multiple layers of features from tiny images.
Tech. Report, University of Toronto (2009)

\bibitem{li2024sigma}
Li, H., et al.:
SIGMA: Sinkhorn-guided masked video modeling.
In: ECCV (2024)

\bibitem{oquab2024dinov2}
Oquab, M., et al.:
DINOv2: Learning robust visual features without supervision.
TMLR (2024)

\bibitem{parkhi2012oxfordpets}
Parkhi, O.M., Vedaldi, A., Zisserman, A., Jawahar, C.V.:
Cats and dogs.
In: CVPR (2012)

\bibitem{peng2022beitv2}
Peng, Z., et al.:
BEiT v2: Masked image modeling with vector-quantized visual tokenizers.
arXiv:2208.06366 (2022)

\bibitem{sablayrolles2019spreading}
Sablayrolles, A., Douze, M., Schmid, C., J\'{e}gou, H.:
Spreading vectors for similarity search.
In: ICLR (2019)

\bibitem{zhou2022ibot}
Zhou, J., et al.:
iBOT: Image BERT pre-training with online tokenizer.
In: ICLR (2022)

\end{thebibliography}

\end{document}
